%-------------------------------------------------------------------------------
%	SECTION TITLE
%-------------------------------------------------------------------------------
\cvsection{Publications}


%-------------------------------------------------------------------------------
%	CONTENT
%-------------------------------------------------------------------------------
\begin{cventries}
\small{

  \begin{cvitems} % Description(s)
      \item \textbf{Yuelong Chen}, Ting Fung Chan. “Initiating the RNA Modification Overture with BALEEN: Reveals the Duets Events in Subcellular Fractions” Target Journal \textbf{\textcolor{red}{Nature Methods}} \textit{(Finalizing)} \textbf{\textcolor{red}{【IF: 36.1】:中科院1区}}
      \item \textbf{Yuelong Chen}, Haonan Wu, Min Cao, Nana Jin, Ning Zhang, Fei Shi, Tao Liu, Kaisheng Liu, Lixin Cheng. “Moonlighting lncRNA PKI55 acts as a molecular switch orchestrating sepsis via module activity reversal” \textbf{\textcolor{red}{IScience}} \textit{(under review)} \textbf{\textcolor{red}{【IF: 4.1】:中科院2区}}
      \item \textbf{Yuelong Chen}, et al. “RNA Modification Modules Assist in Sepsis Diagnosis” \textbf{\textcolor{red}{Briefings in Bioinformatics}} \textit{(under review)} \textbf{\textcolor{red}{【IF: 7.7】:生化研究方法小类1区}}
      \item Nana Jin, Chuanchuan Nan, Wanyang Li, Peijing Lin, Yu Xin, Jun Wang, \textbf{Yuelong Chen}, Yuanhao Wang, Kaijiang Yu, Changsong Wang, Chunbo Chen, Qingshan Geng, Lixin Cheng. “PAGE-based transfer learning from single-cell to bulk sequencing enhances model generalization for sepsis diagnosis.” \textbf{\textcolor{red}{Briefings in Bioinformatics}} 26, no. 1 (January 2025): bbae661. \textbf{\textcolor{red}{【IF: 7.7】:生化研究方法小类1区}}
      \item Jizhou Zhang, Xiao Lin, \textbf{Yuelong Chen}, Tsz-Ho Li, Alan Chun-Kit Lee, Eugene Yui-Ching Chow, William Chi-Shing Cho, and Ting-Fung Chan. “LAFITE Reveals the Complexity of Transcript Isoforms in Subcellular Fractions.” \textbf{\textcolor{red}{Advanced Science}} (Weinheim, Baden-Wurttemberg, Germany) 10, no. 3 (January 2023): e2203480. \textbf{\textcolor{red}{【IF: 14.1】:中科院1区}}
      \item Yanshuo Chen, Yixuan Wang, \textbf{Yuelong Chen}, Yuqi Cheng, Yumeng Wei, Yunxiang Li, Jiuming Wang, Yingying Wei, Ting-Fung Chan, and Yu Li. “Deep Autoencoder for Interpretable Tissue-Adaptive Deconvolution and Cell-Type-Specific Gene Analysis.” \textbf{\textcolor{red}{Nature Communications}} 13, no. 1 (November 8, 2022): 6735. \textbf{\textcolor{red}{【IF: 15.7】:中科院1区}}
      \item Keng Po Lai, Nathan Yi Kan Tam, \textbf{Yuelong Chen}, Chi Tim Leung, Xiao Lin, Chau Fong Tsang, Yin Cheung Kwok, et al. “miRNA–mRNA Integrative Analysis Reveals the Roles of miRNAs in Hypoxia-Altered Embryonic Development- and Sex Determination-Related Genes of Medaka Fish.” \textbf{\textcolor{red}{Frontiers in Marine Science}} 8 (January 21, 2022): 736362. \textbf{\textcolor{red}{【IF: 3.0】:海洋与淡水生物学小类1区}}
      \item Yuanyuan Xu,  Yan Li, Qingyu Xu, \textbf{Yuelong Chen}, Na Lv, Yu Jing, Liping Dou, et al. “Implications of Mutational Spectrum in Myelodysplastic Syndromes Based on Targeted Next-Generation Sequencing.” Oncotarget 8, no. 47 (October 10, 2017): 82475–90. \textbf{\textcolor{red}{【发表时IF: 5.2】:发表时中科院1区}}
      \item Yunfeng Wang, Xiujie Chen, Lei Liu, \textbf{Yuelong Chen}, Hongzhe Ma, Ruizhi Yang, and Xiangqiong Liu. “Identifying the Causative Proteins of Similar Side Effect Pairs to Explore the Common Molecular Basis of These Side Effects.” Molecular bioSystems 11, no. 7 (July 2015): 2060–67. \textbf{\textcolor{red}{【IF: 2.8】}}
      \item Fujian Tan,  Ruizhi Yang, Xiaoxue Xu, Xiujie Chen, Yunfeng Wang, Hongzhe Ma, Xiangqiong Liu,  \textbf{Yuelong Chen}. “Drug Repositioning by Applying ‘expression Profiles’ Generated by Integrating Chemical Structure Similarity and Gene Semantic Similarity.” Molecular bioSystems 10, no. 5 (May 2014): 1126–38. \textbf{\textcolor{red}{【IF: 2.8】}}
      \item Xiujie Chen,  Xiangqiong Liu, Xiaodong Jia, Fujian Tan, Ruizhi Yang, Sheng Chen, Lei Liu, Yunfeng Wang, and \textbf{Yuelong Chen}. “Network Characteristic Analysis of ADR-Related Proteins and Identification of ADR-ADR Associations.” Scientific Reports 3 (2013): 1744. \textbf{\textcolor{red}{【IF: 3.9】:中科院3区}}
  \end{cvitems}
}







%---------------------------------------------------------
\end{cventries}